\chapter[Trabalhos Relacionados]{Trabalhos Relacionados}
\label{cap:relacionados}

Este capítulo situa o trabalho em relação a recursos e estudos de extração de
informação clínica, com ênfase na língua portuguesa.

\section{Recursos e desafios em inglês}

A extração de relações clínicas consolidou-se em língua inglesa em torno de
campanhas de avaliação como as séries i2b2 e n2c2, que disponibilizaram corpora
anotados e tarefas compartilhadas para problemas como relações entre medicamentos,
dosagens e efeitos adversos. Esses esforços estabeleceram protocolos de avaliação
e evidenciaram a importância de tratar adequadamente o desbalanceamento de classes
e o vazamento de dados entre partições — preocupações transpostas para este
trabalho.

\section{Recursos para o português}

O desenvolvimento de PLN clínico em português foi por muito tempo limitado pela
ausência de corpora anotados de acesso controlado. O \textbf{SemClinBr}
\cite{oliveira_2022_semclinbr} preencheu parte dessa lacuna ao disponibilizar mil
documentos clínicos anotados semanticamente, acompanhados de dicionários de
abreviações e de pistas de negação. É o corpus central deste trabalho.

No plano dos modelos, o \textbf{BioBERTpt} \cite{schneider_2020_biobertpt} adaptou
o paradigma de modelos de linguagem contextuais ao domínio clínico e biomédico em
português brasileiro, demonstrando ganhos em reconhecimento de entidades nomeadas
clínicas em relação a modelos de domínio geral. Modelos multilíngues como o
XLM-R \cite{conneau_2020_xlmr} oferecem uma alternativa de base ampla, útil como
ponto de comparação.

\section{Posicionamento deste trabalho}

Enquanto boa parte da literatura em português concentra-se em NER, este trabalho
aborda a etapa subsequente de extração de relações sobre o SemClinBr, com atenção
particular à negação clínica. A contribuição não está em propor uma arquitetura
inédita, mas em estabelecer, de maneira reprodutível e estatisticamente
fundamentada, uma cadeia completa — do texto bruto às métricas com intervalos de
confiança — que sirva de referência comparável para trabalhos futuros sobre o
mesmo corpus.
